% GENERATED FILE. Edit canonical lessons, then run npm run generate:books.
% canonical-source-hash: fnv1a64:4cc317e84126ca42
% canonical-lessons: ES-C696-termometro, ES-C696-pronostico, ES-C696-prevision, ES-C696-humedo, ES-C696-humedad

\chapter{Tiempo y temperatura}
\label{ch:es-tiempo-y-temperatura}
\hlchaptermodality{696}

\begin{chapteropening}
\textbf{By the end of this chapter:} \emph{I can understand a weather forecast and talk about the thermometer and the humidity.}

Complete the last lesson of chapter 696: I can understand a weather forecast and talk about the thermometer and the humidity.
\end{chapteropening}

\section[el termómetro]{el termómetro — the thermometer}

\label{lesson:ES-C696-termometro}



\begin{quote}
Before the new one: say the Spanish for to fish, then the Spanish for the deer.
\end{quote}

\subsection*{You'll want to know: el termómetro}
\textbf{el termómetro} — \textquotedblleft{}the thermometer\textquotedblright{}.

\begin{quote}
El termómetro marca treinta grados.
\end{quote}

\begin{quote}
\emph{The thermometer shows thirty degrees.}
\end{quote}

\emph{Treinta grados a la sombra} — thirty degrees in the shade.

\begin{cousinweb}
\emph{termómetro} came from French \emph{thermomètre}, a word made from Greek \emph{thermós} \textquotedblleft{}hot\textquotedblright{} and \emph{métron} \textquotedblleft{}measure\textquotedblright{}.
\end{cousinweb}

\begin{grammarlens}[title={What you've built}]
One more word for telling, explaining and describing at B1.
\end{grammarlens}

\subsection*{Your turn}
\begin{itemize}
\raggedright
  \item \emph{Say it:} \emph{el termómetro}
  \item \emph{Say it:} \emph{el termómetro}, once more
  \item \emph{Say it:} a sentence of your own with \emph{el termómetro}
  \item \emph{Recall:} say the Spanish for to fish, then the Spanish for the deer, then say \emph{el termómetro} again
\end{itemize}

\begin{tcolorbox}[breakable,colback=teal!4,colframe=teal!35!black,title={Before you move on}]
What does el termómetro mean? (\textquotedblleft{}The thermometer\textquotedblright{}.) Say it once more.
\end{tcolorbox}
\section[el pronóstico]{el pronóstico — the forecast; the prognosis}

\label{lesson:ES-C696-pronostico}



\begin{quote}
Before the new one: say the Spanish for the deer, then the Spanish for the thermometer.
\end{quote}

\subsection*{You'll want to know: el pronóstico}
\textbf{el pronóstico} — \textquotedblleft{}the forecast; the prognosis\textquotedblright{}.

\begin{quote}
El pronóstico dice que va a llover.
\end{quote}

\begin{quote}
\emph{The forecast says it is going to rain.}
\end{quote}

\emph{El pronóstico del tiempo} — the weather forecast.

\begin{cousinweb}
\emph{pronóstico} is a learned borrowing from Late Latin \emph{prognōsticus}, from Greek \emph{prognōstikós} \textquotedblleft{}knowing beforehand\textquotedblright{}.
\end{cousinweb}

\begin{grammarlens}[title={What you've built}]
One more word for telling, explaining and describing at B1.
\end{grammarlens}

\subsection*{Your turn}
\begin{itemize}
\raggedright
  \item \emph{Say it:} \emph{el pronóstico}
  \item \emph{Say it:} \emph{el pronóstico}, once more
  \item \emph{Say it:} a sentence of your own with \emph{el pronóstico}
  \item \emph{Recall:} say the Spanish for the deer, then the Spanish for the thermometer, then say \emph{el pronóstico} again
\end{itemize}

\begin{tcolorbox}[breakable,colback=teal!4,colframe=teal!35!black,title={Before you move on}]
What does el pronóstico mean? (\textquotedblleft{}The forecast; the prognosis\textquotedblright{}.) Say it once more.
\end{tcolorbox}
\section[la previsión]{la previsión — the outlook, the forecast; foresight}

\label{lesson:ES-C696-prevision}



\begin{quote}
Before the new one: say the Spanish for the thermometer, then the Spanish for the forecast.
\end{quote}

\subsection*{You'll want to know: la previsión}
\textbf{la previsión} — \textquotedblleft{}the outlook, the forecast; foresight\textquotedblright{}.

\begin{quote}
La previsión para mañana es buena.
\end{quote}

\begin{quote}
\emph{The forecast for tomorrow is good.}
\end{quote}

The verb is \emph{prever}.

\begin{cousinweb}
\emph{previsión} is a learned borrowing from Late Latin \emph{praevīsiō} \textquotedblleft{}foresight\textquotedblright{} (EtymDB), from \emph{praevidēre} \textquotedblleft{}to foresee\textquotedblright{}, built on \emph{vidēre} \textquotedblleft{}to see\textquotedblright{}.
\end{cousinweb}

\begin{grammarlens}[title={What you've built}]
One more word for telling, explaining and describing at B1.
\end{grammarlens}

\subsection*{Your turn}
\begin{itemize}
\raggedright
  \item \emph{Say it:} \emph{la previsión}
  \item \emph{Say it:} \emph{la previsión}, once more
  \item \emph{Say it:} a sentence of your own with \emph{la previsión}
  \item \emph{Recall:} say the Spanish for the thermometer, then the Spanish for the forecast, then say \emph{la previsión} again
\end{itemize}

\begin{tcolorbox}[breakable,colback=teal!4,colframe=teal!35!black,title={Before you move on}]
What does la previsión mean? (\textquotedblleft{}The outlook, the forecast; foresight\textquotedblright{}.) Say it once more.
\end{tcolorbox}
\section[húmedo]{húmedo — humid, damp}

\label{lesson:ES-C696-humedo}



\begin{quote}
Before the new one: say the Spanish for the forecast, then the Spanish for the outlook.
\end{quote}

\subsection*{You'll want to know: húmedo}
\textbf{húmedo} — \textquotedblleft{}humid, damp\textquotedblright{}.

\begin{quote}
El aire es muy húmedo en la costa.
\end{quote}

\begin{quote}
\emph{The air is very humid on the coast.}
\end{quote}

\emph{Ropa húmeda} — damp clothes.

\begin{cousinweb}
\emph{húmedo} is a learned borrowing from Latin \emph{hūmidus} \textquotedblleft{}moist, damp\textquotedblright{} (EtymDB).
\end{cousinweb}

\begin{grammarlens}[title={What you've built}]
One more word for telling, explaining and describing at B1.
\end{grammarlens}

\subsection*{Your turn}
\begin{itemize}
\raggedright
  \item \emph{Say it:} \emph{húmedo}
  \item \emph{Say it:} \emph{húmedo}, once more
  \item \emph{Say it:} a sentence of your own with \emph{húmedo}
  \item \emph{Recall:} say the Spanish for the forecast, then the Spanish for the outlook, then say \emph{húmedo} again
\end{itemize}

\begin{tcolorbox}[breakable,colback=teal!4,colframe=teal!35!black,title={Before you move on}]
What does húmedo mean? (\textquotedblleft{}Humid, damp\textquotedblright{}.) Say it once more.
\end{tcolorbox}
\section[la humedad]{la humedad — the humidity, the damp}

\label{lesson:ES-C696-humedad}



\begin{quote}
Before the new one: say the Spanish for the outlook, then the Spanish for humid.
\end{quote}

\subsection*{You'll want to know: la humedad}
\textbf{la humedad} — \textquotedblleft{}the humidity, the damp\textquotedblright{}.

\begin{quote}
Con tanta humedad hace más calor.
\end{quote}

\begin{quote}
\emph{With so much humidity it feels hotter.}
\end{quote}

\emph{Hay humedad en la pared} — there is damp in the wall.

\begin{cousinweb}
\emph{humedad} is built on \emph{húmedo} \textquotedblleft{}humid\textquotedblright{}, a learned borrowing from Latin \emph{hūmidus} \textquotedblleft{}moist\textquotedblright{}.
\end{cousinweb}

\begin{grammarlens}[title={What you've built}]
One more word for telling, explaining and describing at B1.
\end{grammarlens}

\subsection*{Your turn}
\begin{itemize}
\raggedright
  \item \emph{Say it:} \emph{la humedad}
  \item \emph{Say it:} \emph{la humedad}, once more
  \item \emph{Say it:} a sentence of your own with \emph{la humedad}
  \item \emph{Recall:} say the Spanish for the outlook, then the Spanish for humid, then say \emph{la humedad} again
\end{itemize}

\begin{tcolorbox}[breakable,colback=teal!4,colframe=teal!35!black,title={Before you move on}]
What does la humedad mean? (\textquotedblleft{}The humidity, the damp\textquotedblright{}.) Say it once more.
\end{tcolorbox}
